\chapter{Desarrollo y Consignas}

\section{Implementación del Amplificador}
Implementar un amplificador realimentado según la topología del circuito de la Figura 1.

\vspace{1cm}
% Espacio para incluir la imagen del esquemático o dibujarlo con circuitikz
\begin{figure}[H]
    \centering
    % \includegraphics[width=0.8\textwidth]{images/esquematico.png}
    % \caption{Circuito esquemático del Amplificador de 2 etapas.}
    % \label{fig:esquematico}
    [AQUÍ VA EL ESQUEMÁTICO DEL CIRCUITO - COMPLETAR CON IMAGEN O TIKZ]
\end{figure}

\subsection{Valores de Componentes Normalizados}
A continuación se detallan los valores de los componentes a utilizar para el armado del circuito:

\begin{table}[H]
    \centering
    \begin{tabular}{|l|l||l|l||l|l||l|l|}
        \hline
        $V_{cc}$ & 22 $\V$ & $R_{e1}$ & 130 $\ohm$ & $R_{c2}$ & 2.0 $\kohm$ & $R_f$ & 1.6 $\kohm$ \\ \hline
        $R_{b1}$ & 820 $\kohm$ & $R_{e2}$ & 12 $\kohm$ & $R_{e3}$ & 240 $\ohm$ & $C_{a1}$ & 1 $\uF$ \\ \hline
        $R_{b2}$ & 560 $\kohm$ & $R_{b3}$ & 110 $\kohm$ & $R_{e4}$ & 200 $\ohm$ & $C_{E1}$ & 10 $\uF$ \\ \hline
        $R_{c1}$ & 18 $\kohm$ & $R_{b4}$ & 24 $\kohm$ & $R_L$ & 470 $\ohm$ & $C_{a2}$ & 1 $\uF$ \\ \hline
        $C_{e2}$ & 10 $\uF$ & $C_o$ & 1 $\uF$ & $Q_1$ & BC 337 & $Q_2$ & BC 337 \\ \hline
    \end{tabular}
    \caption{Valores de Componentes Normalizados}
    \label{tab:componentes}
\end{table}

\section{Análisis Teórico y Cálculos}
[COMPLETAR CON LOS CÁLCULOS TEÓRICOS DEL PUNTO DE OPERACIÓN Y PARÁMETROS DINÁMICOS]
